% ---- p2-namespaced TikZ styles (package loading happens in preamble.txt) ----
\tikzset{
  p2ent/.style={circle, draw=black!65, line width=0.7pt, minimum size=9.5mm,
                inner sep=0pt, font=\small\bfseries},
  p2lab/.style={font=\scriptsize},
  p2rel/.style={draw=black!65, line width=0.7pt},
  p2el/.style={font=\scriptsize, fill=white, inner sep=1.5pt},
}
\newcommand{\pTwoTgt}{orange!22}
\newcommand{\pTwoAgt}{blue!12}
\newcommand{\pTwoAtt}{green!14}
\newcommand{\pTwoVen}{purple!12}

\subsection*{1. Set-up: HINs, network schemas and meta-paths}

\paragraph{Heterogeneous information network (HIN).}
Following Sun et al.\ (\emph{PathSim}, PVLDB 4(11):992--1003, 2011) and the survey of Shi, Li,
Zhang, Sun \& Yu (\emph{A Survey of Heterogeneous Information Network Analysis}, IEEE TKDE
29(1):17--37, 2017), a \textbf{heterogeneous information network} is a graph
\[
  G=(V,E,\phi,\psi),\qquad \phi:V\to\mathcal{A},\qquad \psi:E\to\mathcal{R},
\]
in which every node $v\in V$ carries exactly one \emph{node type} $\phi(v)$ from the type set
$\mathcal{A}$, every edge $e\in E$ carries exactly one \emph{relation type} $\psi(e)$ from
$\mathcal{R}$, and
\[
  \boxed{\;|\mathcal{A}|+|\mathcal{R}|>2\;}
\]
This inequality \emph{is} the definition of heterogeneity: when $|\mathcal{A}|=|\mathcal{R}|=1$
the network degenerates to the homogeneous graph of Problem~1. Each relation is typed at both
ends, inducing a map $A_i\xrightarrow{R}A_j$ between two specific node types; its inverse
$R^{-1}$ is in general a \emph{different} relation ($\mathsf{writes}$ vs.\ $\mathsf{written\text{-}by}$),
which is why the released benchmarks below count $6$ or $8$ ``edge types'' where a reader might
expect $3$ or $4$.

\paragraph{Network schema.}
The \textbf{network schema} $T_G=(\mathcal{A},\mathcal{R})$ is the directed graph whose vertices
are the node types $\mathcal{A}$ and whose edges are the relation types $\mathcal{R}$. It is a
meta-template of which $G$ is an instantiation: every node and edge of $G$ maps to some type in
$T_G$, so the schema states which kinds of connection the dataset can express at all. This is the
object Section~\ref{p2:sec:schemas} draws for each benchmark. Note that the drawing carries real
information here whereas it does not in the homogeneous case: Cora's schema is a single node type
with one self-relation, while DBLP's is a four-type star.

\paragraph{Meta-path.}
A \textbf{meta-path} $\mathcal{P}$ of length $\ell$ is a path defined \emph{on the schema},
\[
  \mathcal{P}\;:\;A_1\xrightarrow{R_1}A_2\xrightarrow{R_2}\cdots\xrightarrow{R_\ell}A_{\ell+1},
\]
denoting the composite relation $R=R_1\circ R_2\circ\cdots\circ R_\ell$ from type $A_1$ to type
$A_{\ell+1}$. A node sequence $v_1v_2\cdots v_{\ell+1}$ in $G$ with $\phi(v_i)=A_i$ and
$\psi(v_iv_{i+1})=R_i$ is a \emph{path instance}. Meta-paths are written by type initials, so on
DBLP $\mathrm{APA}$ abbreviates
$\text{Author}\xrightarrow{\mathsf{writes}}\text{Paper}\xrightarrow{\mathsf{written\text{-}by}}\text{Author}$.

Meta-paths matter because each defines a \emph{different} homogeneous projection of the same HIN,
and the projections carry different semantics: on DBLP $\mathrm{APA}$ means \emph{co-authorship}
while $\mathrm{APCPA}$ means \emph{published at the same venue}, and two authors may be strongly
related under one and unrelated under the other. Nearly every method in
Section~\ref{p2:sec:baselines} is a different answer to ``how should meta-path structure be
used?'' --- fixed and hand-chosen (metapath2vec, HAN), aggregated with attention (MAGNN), learned
(GTN), or discarded in favour of per-relation parameters (R-GCN, HGT, simple-HGN).

% =====================================================================
